\documentclass[10pt]{article}

\usepackage[utf8]{inputenc}

\usepackage[T1]{fontenc}

\usepackage{amsmath}

\usepackage{amsfonts}

\usepackage{amssymb}

\usepackage[version=4]{mhchem}

\usepackage{stmaryrd}

\usepackage{graphicx}

\usepackage[export]{adjustbox}

\graphicspath{ {./images/} }



\begin{document}

выми множествами. Так, в 1834 Н. И. Лобачевский [2] писал: «Общее шонятие функции требует, чтобы функцией от $x$ называть число, которое дается для каждого $x$ и вместе с $x$ постепенно изменяется. Значение функции может быть дано или аналитическим выражением, или условием, к-рое подает средство испытывать все числа и выбрать одно из них; или, наконец, зависимость может существовать и оставаться неизвестной». Определение Ф. как соответствия между двумя произвольными (не обязательно числовыми) множествами в 1887 было сформули ровано Р. Дедекиндом (R. Dedekind) [3].



Понятие соответствия, а следовательно, и понятие Ф., иногда сводится к другим понятиям (множеству [4], отношению [5] или несколько другим теоретико-множественным и логико-математич. концепциям [6]), а иногда принимается за первичное, неопределимое понятие [7], поскольку, как әто выразил, напр., А. Чёрч (A. Church): «В конечном счете понятие функции или какое-либо сходное понятие, напр., понятие класса,- приходится считать первоначальным, или неопределимым» [8].



Ниже рассматривается понятие Ф., основанное на понятии множества и простейших операций над множествами.



Говорят, что число элементов множества $\boldsymbol{A}$ равно 1, или что множество $A$ состоит из одного әлемента, если в нем имеется әлемент $a$ и нет других (иначе говоря, если после вычитания из множества $A$ множества $\{a\}$ получается пусгое множество). Непустое множество $A$ наз. множеством из двух әлементов или п а р о й: $A=\{a, b\}$, если после вычитания из него множества, состоящего только из одного әлемента $a \in A$, останется множество, к-рое также состоит из одного элемента $b \in A$ (это определение не зависит от выбора указанного әлемента $a \in A)$.



Если задана пара $A=\{a, b\}$, то пара $\{a,\{a, b\}\}$ наз. у по р я доч е н н о й п а р о й әлементов $a \in A$ и $b \in A$ и обозначается $(a, b)$. Элемент $a \in A$ наз. ее шервым элементом, а элемент $b \in A$ - вторым.



Если заданы множества $X$ и $\bar{Y}$, то множество всех упорядоченных пар $(x, y), x \in X, y \in Y$, наз. п $\mathrm{p} о$ и зв е д ен и е м м н о ж е с т в $X$ и $Y$ и обозначается через $X \times Y$. При әтом не предполагается, что обязательно множество $X$ отлично от множества $Y$, т. е. возможен случай, когда $X=Y$.



Всякое множество $f=\{(x, y)\}$ упорядоченных пар $(x, y), x \in X, y \in Y$, такое, что для любых пар $\left(x^{\prime}, y^{\prime}\right) \in f$ и $\left(x^{\prime \prime}, y^{\prime \prime}\right) \in f$ из условия $y^{\prime} \neq y^{\prime \prime}$ следует, что $x^{\prime} \neq x^{\prime \prime}$, наз. ф у нк ци е й, или, что то же самое, о т о б р аже ни е м.



Наряду с терминами «Ф.» и «отображение» употребляются равнозначные им в определенных ситуациях термины «преобразование», «морфизм», «соответствие».



Множество всех первых элементов упорядоченных пар $(x, y)$ нек-рой функции $f$ наз. областью определения (или множеством определения) этой функции и обозначается $X_{f}$, а множество всех вторых элементов областью зиачений (множеством значений) и обозначается $Y_{f}$. Само множество упорядоченных пар $f=$ $=\{(x, y)\}$, рассматриваемое как подмножество произведения $X \times Y$, наз. г р а ф и к о м функции $f$.



Элемент $x \in X_{f}$ наз. а р г у м ен т о м Ф. или н е-



\includegraphics[max width=\textwidth, center]{2023_07_09_31756bd82f14506c85d4g-1}

з а в и и мо й п е p е мен но й.



Если $f=\{(x, y)\}$ есть функция, то пишут $f: X_{f} \rightarrow Y$ и говорят, что $f$ отображает множество $X_{f}$ в множество $Y$. В случае $X=X_{f}$ пипется просто $f: X \rightarrow Y$.



Если $f: X \longrightarrow Y-\Phi$. и $(x, y) \in f$, то пишут $y=f(x)$ (иногда просто $y=f x$, или $y=x f$ ), а также $f: x \longmapsto y$, $x \in X, y \in Y$, и говорят, что Ф. $f$ ставит в соответствие әлементу $x$ әлемен' $y$ (отображение $f$ отображает әлемент $x$ в элемент $y$ ) или, тто тоже самое, әлемент $y$ соответствует әлементу $x$. В этом с лучае говорят также, что элемент $y$ является значением Ф. $f$ в точке $x$, или о б р азом элем ента $x$ при отображении $f$.



Наряду с символом $f\left(x_{0}\right)$ для обозначенія значения $\Phi . f$ в точке $x_{0}$ употребляется также обозначение $\left.f(x)\right|_{x=x_{0}}$.



Иногда сама функция $f$ обозначается символом $f(x)$. Обозначение $\Phi . f: X \longrightarrow Y$ и ее значения в точке $x \in X$ одним и тем же символом $f(x)$ обычно це приводит к недоразумению, так как в каждом конкретном случае, как правило, всегда бывает ясно, о чем именно идет речь. Обозначение $f(x)$ тасто оказывается удобнее обозначения $f: x \longmapsto y$ при вычислениях. Напр., запись $f(x)=x^{2}$ удобнее и проще использовать при аналитич. преобразованиях, чем запись $f: x \longmapsto x^{2}$.



При заданном $y \in Y$ совокупность всех таких элементов $x \in X$, что $f(x)=y$ наз. I р о о б р а 3 о м э л еме н т а $y$ и обозначается $f^{-1}(y)$. Таким образом,



\[

f^{-1}(y)=\{x: x \in X, f(x)=y\}

\]



Очевидно, если $y \in Y \backslash Y_{f}$, то $f^{-1}(y)=\varnothing$.



Пусть задано отображение $f: X \longrightarrow Y$. Иначе говоря, каждому әлементу $x \in X$ поставлен в соответствие і притом единственный әлемент $y \in Y$, и каждый әлемент $y \in Y_{f} \subset Y$ поставлен в соответствие хотя бы одному элементу $x \in X$. Если $Y=X$, то говорят, что отображение $f$ отображает множество $X$ в с е б я. Если $Y=Y_{f}$, т. е. множество $Y$ совпадает с множеством значений функции $f$, то говорят, тто $f$ отображает множество $X$ н а мно ж е с т в о $Y$, или тто отображение $f$ является с ю р в е к т в вым о тоб р а жен и е м, короче сюръекцией. Таким образом, отображение $f: X \longrightarrow Y$ есть сюръекция, если для любого элемента $y \in Y$ существует, по крайней мере, один такой әлемент $x \in X$, что $f(x)=y$.



Если при отображении $f: X \longrightarrow Y$ разным элементам $x \in X$ соответствуют разные элементы $y \in Y$, т. е. при $x^{\prime} \neq x^{\prime \prime}$ имеет место $f\left(x^{\prime}\right) \neq f\left(x^{\prime \prime}\right)$, то отображение $f$ наз. в в $Y$, а также о днол и с тным отоб г а же и е м, или инъекцией. Таким образом, отображение $f: X \longrightarrow Y$ однолистно (инъективно) тогда и только тогда, когда прообраз каждого элемента $y$ принадлежащего множеству значений функции $f$, т. е. $y \in Y_{f}$, состоит в точности из одного элемента. Если отображение $f: X \longrightarrow Y$ является одновременно взаимно однозначным и на множество $Y$ (см. Взаимно однозначное соответствие), т. е. является одновременно инъекцией и сюръекцией, то оно наз. би е к т и вным о тоб р а жени ем, или биекцией.



Если $f: X \rightarrow Y$ и $A \subset X$, то мнонество



\[

S=\{y: y \in Y, y=f(x), x \in A\},

\]



т. е. множество всех тех $y$, в каждый из к-рых при отображении $f$ отображается хотя бы один элемент из подмножества $A$ множества $X$, наз. о б р а 3 о м по дм н о ж е с т в а $A$ и пишется $S=f(A)$. В частности, всегда $Y_{f}=f(X)$. Для образов множеств $A \subset X$ и $B \subset X$ справедливы следующие соотношения



\[

\begin{gathered}

f(A \cup B)=f(A) \cup f(B), \\

f(A \cap B) \subset f(A) \cap f(B), \\

f(A) \backslash f(B) \subset f(A \backslash B),

\end{gathered}

\]



а если $A \subset B$, то $f(A) \subset f(B)$.



Если $f: X \rightarrow Y$ и $S \subset Y$, то множество



\[

A=\{x: x \in X, f(x) \in S\}

\]



наз. прооб разом множества $S$ и пишут $A=f^{-1}(S)$. Таким образом, прообраз множества $S$ состоит из всех тех әлементов $x \in X$, к-рые при отображении $f$ отображаются в әлементы из $S$, или, что





\end{document}